\chapter{Architecture of the program}
\label{ch:architecture}

Texture Explorer is about 1\,800 lines of C++20 and HLSL. It uses three small libraries:
Dear ImGui for the user interface, \code{stb\_image} to decode images, and \code{nlohmann/json}
to read the list of materials. All three are fetched by CMake when the project is configured.

\section{Modules}

\begin{table}[ht]
\centering\small
\begin{tabularx}{\linewidth}{@{}l>{\raggedright\arraybackslash}X@{}}
\toprule
\textbf{Files} & \textbf{Responsibility}\\
\midrule
\file{src/Mesh.h}, \file{.cpp} & Parametric patches of the four solids; mesh generation; tangent frames; inverting the texture mapping\\
\file{src/Material.h}, \file{.cpp} & The manifest of materials; loading images into GPU textures and CPU copies; bilinear sampling on the CPU\\
\file{src/Inspector.h}, \file{.cpp} & The probe: map values, slope, tangent-space and world-space normals at a texel\\
\file{shaders/pbr.hlsl} & Displacement (vertex shader); bump mapping and GGX lighting (pixel shader); debug views\\
\file{shaders/gizmo.hlsl} & Coloured lines for the arrows of the probe and the light marker\\
\file{src/Renderer.h}, \file{.cpp} & Direct3D~11 device, swap chain, shader compilation, off-screen multisampled target, draw calls\\
\file{src/App.h}, \file{.cpp} & Application state and the six panels of the user interface\\
\file{src/main.cpp} & Window, Dear ImGui set-up, main loop\\
\file{tools/fetch\_textures.ps1} & Downloads the ten materials and writes \file{assets/materials/manifest.json}\\
\bottomrule
\end{tabularx}
\caption{The source files of Texture Explorer.}
\label{tab:modules}
\end{table}

\section{One frame}

\begin{figure}[ht]
\centering
\begin{tikzpicture}[node distance=7mm and 9mm,font=\small,
  box/.style={draw,rounded corners=2pt,fill=#1,align=center,minimum height=9mm,text width=31mm},
  box/.default=black!4, >={Stealth}]
  \node[box=black!8] (ui) {\textbf{App::Frame}\\build the panels};
  \node[box,right=of ui] (maps) {Maps / Inspector\\images: mouse $\to(u,v)$};
  \node[box,right=of maps] (insp) {\textbf{Inspect}\\(CPU, one texel)};
  \node[box,below=of insp] (lines) {arrows $\Nv,\vect n',\Tv,\Bv$\\light marker};
  \node[box,below=of maps] (fc) {\code{FrameConstants}\\camera, light, settings};
  \node[box=npcolor!12,below=of fc] (gpu) {\textbf{RenderScene} (GPU)\\VS: displacement\\PS: bump + GGX};
  \node[box,left=of gpu] (img) {off-screen image\\shown in Viewport};
  \node[box=black!8,above=of maps] (lib) {\textbf{MaterialLibrary}\\CPU copies + textures};
  \draw[->] (ui) -- (maps);
  \draw[->] (maps) -- node[above]{probe} (insp);
  \draw[->] (insp) -- (lines);
  \draw[->] (lines) |- (gpu);
  \draw[->] (fc) -- (gpu);
  \draw[->] (gpu) -- (img);
  \draw[->,dashed] (lib) -| (ui);
  \draw[->,dashed] (lib) -| (insp);
  \draw[->,dashed] (lib.east) -- (lib.east -| insp.east) -- ++(0.6,0)
     |- ($(gpu.south)+(0,-0.5)$) -- (gpu.south);
\end{tikzpicture}
\caption{The flow of one frame. The same material data feeds the probe (on the CPU) and the
renderer (on the GPU), which apply the same formulas.}
\label{fig:frame}
\end{figure}

Dear ImGui is an \emph{immediate-mode} library: the user interface is not a tree of widget
objects but a function that runs every frame and calls \code{ImGui::Button},
\code{ImGui::SliderFloat}, and so on. Each call draws a widget and tells whether the user
interacted with it. The application's state therefore lives in ordinary member variables that
the rest of the program reads directly. One frame proceeds as follows
(Figure~\ref{fig:frame}):
\begin{enumerate}
\item The animations advance by the elapsed time.
\item The panels are built. While the \emph{Maps} and \emph{Inspector} panels draw their images,
they convert the mouse position into texture coordinates and update the \emph{probe}.
\item If the probe is valid, \code{Inspect} computes everything about that texel, on the CPU.
\item The \emph{Viewport} panel fills a constant buffer with the camera, the light and the
settings, adds the inspector's arrows to a list of lines, and asks the renderer to draw the scene
into an off-screen image of the size of the panel. It then shows that image.
\item Dear ImGui draws the whole interface into the window, and the frame is presented.
\end{enumerate}

\section{A design principle: compute it twice}
The pixel shader computes the bumped normal for millions of pixels and shows only the result.
To make the computation visible, the probe repeats it on the CPU for one texel and shows every
intermediate value. The two implementations must agree exactly, so they are written to be
identical:
\begin{itemize}
\item the CPU keeps a byte-for-byte copy of every map and samples it bilinearly with the GPU's
texel-centre convention (Section~\ref{sec:cpusampling});
\item the slope is estimated with the same central differences, at level~0, with the same
strength (\code{BumpNormalTangent} in C++, \code{BumpNormalTS} in HLSL);
\item the tangent frame is built with the same Gram--Schmidt construction from the same patch
functions (\code{TangentFrame} and the pixel shader);
\item the 3D position of the probe comes from the same parametric functions that generated the
mesh, so the arrows land on the pixels that show the probed texel.
\end{itemize}
The probe and the image disagree only where the program simplifies on purpose. For example, the
probe evaluates the exact surface, while the GPU draws a mesh of flat triangles.
